\chapter{Line and Surface Integrals}
\label{ch:linesurface}

So far we have integrated over flat regions. Many quantities, however, are defined along curves (work, circulation) or across surfaces (flux, surface charge). In this chapter we define line integrals and surface integrals, prove that they do not depend on how the curve or surface is parametrized, and prove the fundamental theorem for line integrals. We also study the central example that shows why topology matters: the vortex field.

%====================================================================
\section{Line integrals of scalar fields}
%====================================================================

Let $C$ be a curve with a regular $C^1$ parametrization $\vec r:[a,b]\to\R^3$ (see Chapter~\ref{ch:geometry}). The \emph{line integral of a scalar function $f$ with respect to arc length} is
\[
\int_Cf\,\dd s=\int_a^bf\bigl(\vec r(t)\bigr)\,\norm{\vec r\,'(t)}\,\dd t .
\]
Here $\dd s=\norm{\vec r\,'(t)}\,\dd t$ is the arc-length element. For $f\equiv1$ we get the length of $C$. If $f$ is a density (mass per unit length) along a wire, $\int_Cf\,\dd s$ is the total mass of the wire.

\begin{proposition}[Independence of parametrization]\label{prop:lineparam}
If $\vec\rho(s)=\vec r(\varphi(s))$ where $\varphi:[c,d]\to[a,b]$ is a $C^1$ bijection (increasing or decreasing), then $\int_c^df(\vec\rho(s))\norm{\vec\rho\,'(s)}\dd s=\int_a^bf(\vec r(t))\norm{\vec r\,'(t)}\dd t$.
\end{proposition}

\begin{proof}
$\norm{\vec\rho\,'(s)}=\norm{\vec r\,'(\varphi(s))}\,\abs{\varphi'(s)}$. Substitute $t=\varphi(s)$. If $\varphi'>0$ the limits are $a,b$ and $\dd t=\varphi'\dd s$. If $\varphi'<0$ the limits are reversed, which gives a factor $-1$ that cancels the sign in $\abs{\varphi'}=-\varphi'$. In both cases we obtain the right-hand side.
\end{proof}

So the scalar line integral depends on the \emph{curve as a set}, not on its parametrization or direction.

%====================================================================
\section{Line integrals of vector fields: work and circulation}
%====================================================================

Let $\vec F$ be a continuous vector field and $C$ an \emph{oriented} curve (a curve with a chosen direction of travel), parametrized by $\vec r(t)$ in the direction of travel. The \emph{line integral of $\vec F$ along $C$} is
\[
\int_C\vec F\cdot\dd\vec r=\int_a^b\vec F\bigl(\vec r(t)\bigr)\cdot\vec r\,'(t)\,\dd t=\int_C\vec F\cdot\vec T\,\dd s,
\]
where $\vec T=\vec r\,'/\norm{\vec r\,'}$ is the unit tangent. It adds up the component of $\vec F$ \emph{along} the curve. If $\vec F$ is a force, this is the \emph{work} done along $C$. If $C$ is closed and $\vec F$ is a velocity field, it is the \emph{circulation}, and we write $\oint_C\vec F\cdot\dd\vec r$. In coordinates, $\vec F\cdot\dd\vec r=P\,\dd x+Q\,\dd y+R\,\dd z$.

\begin{figure}[ht]
\centering
\begin{tikzpicture}[scale=1.35]
  % curve
  \draw[curvestyle,thick] (0,0) to[out=60,in=200] (2,1.7) to[out=20,in=190] (4.6,1.4);
  \node[accent,font=\small] at (4.6,1.0) {$C$};
  % point and vectors at curve point P ~ (2,1.7)
  \coordinate (P) at (2,1.7);
  \node[pt] at (P) {};
  % tangent direction approx (1,0.15)->(0.99,0.15)
  \draw[vecstyle,accent3] (P)--++(0.99,0.15) node[above right,font=\small]{$\vec T$};
  \draw[vecstyle,accent2] (P)--++(0.9,0.9) node[above right,font=\small]{$\vec F$};
  \draw[guide] (P)++(0.9,0.9)--++(0.09,-0.75);
  \draw[very thick,accent4] (P)--++(0.9,0.135);
  \node[accent4,font=\small] at (3.0,1.65) [below] {$\vec F\cdot\vec T$};
  % several small arrows along the curve to indicate dr
  \foreach \t in {0.15,0.35,0.55}{
    \pgfmathsetmacro{\xx}{4.6*\t*0.85}
  }
  \node[font=\small] at (0.1,-0.35) {start};
  \node[font=\small] at (4.9,1.75) {end};
\end{tikzpicture}
\caption{A line integral adds up the component of $\vec F$ along the direction of travel. At each point only the projection $\vec F\cdot\vec T$ counts; the component of $\vec F$ perpendicular to the curve does no work.}
\label{fig:lineint}
\end{figure}

\begin{proposition}[Properties of the line integral]\label{prop:lineprops}
\begin{enumerate}[label=(\alph*)]
  \item The integral is linear in $\vec F$ and additive over a subdivision $C=C_1+C_2$.
  \item Reversing the orientation changes the sign: $\int_{-C}\vec F\cdot\dd\vec r=-\int_C\vec F\cdot\dd\vec r$.
  \item If $\vec\rho=\vec r\circ\varphi$ with $\varphi'>0$, then $\vec\rho$ and $\vec r$ give the same integral. So the integral depends only on the oriented curve.
\end{enumerate}
\end{proposition}

\begin{proof}
(a) is clear from the definition. (b) and (c): let $\vec\rho(s)=\vec r(\varphi(s))$. Then $\vec F(\vec\rho)\cdot\vec\rho\,'(s)=\vec F(\vec r(\varphi))\cdot\vec r\,'(\varphi)\,\varphi'(s)$, and the substitution $t=\varphi(s)$ gives $\int_c^d\bigl[\vec F(\vec r(t))\cdot\vec r\,'(t)\bigr]_{t=\varphi(s)}\varphi'(s)\dd s=\int_{\varphi(c)}^{\varphi(d)}\vec F(\vec r(t))\cdot\vec r\,'(t)\dd t$. If $\varphi'>0$ then $\varphi(c)=a$, $\varphi(d)=b$ and we get (c). If $\varphi'<0$ the limits are $b$ and $a$, and the result is the negative: this is (b) (the curve $\vec\rho$ is $\vec r$ traversed backwards).
\end{proof}

\begin{example}[Circulation of the rotation field]
Let $\vec F=(-y,x)$ and let $C$ be the circle of radius $R$, counterclockwise: $\vec r(t)=(R\cos t,R\sin t)$, $0\le t\le2\pi$. Then $\vec F(\vec r(t))=(-R\sin t,R\cos t)=\vec r\,'(t)$ and
\[
\oint_C\vec F\cdot\dd\vec r=\int_0^{2\pi}R^2\,\dd t=2\pi R^2=2\cdot\Area(\text{disc}).
\]
This matches the curl: $\curl\vec F=2\khat$. We will see in Chapter~\ref{ch:theorems} (Green's theorem) that this is no coincidence.
\end{example}

\begin{example}[Work along different paths]
Let $\vec F=(y,\,x)$ and compare two paths from $(0,0)$ to $(1,1)$. Path 1: $\vec r=(t,t)$, $\vec F\cdot\vec r\,'=(t,t)\cdot(1,1)=2t$, integral $=1$. Path 2: first along the $x$-axis to $(1,0)$ ($\vec F=(0,x)$, $\vec r\,'=(1,0)$: integral $0$), then up to $(1,1)$ ($\vec F=(y,1)$, $\vec r\,'=(0,1)$: integral $\int_0^1 1\,\dd y=1$). Both give $1$. The field is $\nabla(xy)$, so this is expected, as we prove next.
\end{example}

%====================================================================
\section{The fundamental theorem for line integrals}
%====================================================================

\begin{theorem}[Fundamental theorem for line integrals]\label{thm:ftli}
Let $f$ be $C^1$ on an open set containing the oriented curve $C$ with parametrization $\vec r:[a,b]\to\R^n$ (piecewise $C^1$). Then
\[
\int_C\nabla f\cdot\dd\vec r=f\bigl(\vec r(b)\bigr)-f\bigl(\vec r(a)\bigr).
\]
\end{theorem}

\begin{proof}
By the chain rule, $\dfrac{\dd}{\dd t}f(\vec r(t))=\nabla f(\vec r(t))\cdot\vec r\,'(t)$. Integrating from $a$ to $b$ and using the fundamental theorem of calculus,
\[
\int_a^b\nabla f(\vec r(t))\cdot\vec r\,'(t)\,\dd t=\int_a^b\frac{\dd}{\dd t}f(\vec r(t))\,\dd t=f(\vec r(b))-f(\vec r(a)).
\]
For a piecewise $C^1$ curve, apply this to each piece; the sum telescopes.
\end{proof}

\begin{theorem}[Path independence]\label{thm:pathind}
Let $U\subseteq\R^n$ be open and connected and let $\vec F:U\to\R^n$ be continuous. The following are equivalent:
\begin{enumerate}[label=(\alph*)]
  \item $\vec F$ is conservative: $\vec F=\nabla f$ for some $C^1$ function $f$ on $U$.
  \item The line integral $\int_C\vec F\cdot\dd\vec r$ depends only on the endpoints of $C$ (for piecewise $C^1$ curves $C$ in $U$).
  \item $\oint_C\vec F\cdot\dd\vec r=0$ for every closed piecewise $C^1$ curve $C$ in $U$.
\end{enumerate}
\end{theorem}

\begin{proof}
(a)$\Rightarrow$(b): Theorem~\ref{thm:ftli}.

(b)$\Leftrightarrow$(c): Let $C_1,C_2$ be two paths from $\vec p$ to $\vec q$. The closed curve $C_1+(-C_2)$ has integral $\int_{C_1}-\int_{C_2}$ (Proposition~\ref{prop:lineprops}). So the two integrals agree for all such pairs if and only if all closed-curve integrals vanish. (Conversely, a closed curve passes through two points and can be viewed as two paths between them.)

(b)$\Rightarrow$(a): Fix $\vec x_0\in U$. Since $U$ is open and connected, every point $\vec x\in U$ can be joined to $\vec x_0$ by a piecewise linear path. Define $f(\vec x)=\int_{\vec x_0}^{\vec x}\vec F\cdot\dd\vec r$ along any such path; by (b) this is well defined. For small $h$, the segment from $\vec x$ to $\vec x+h\vec e_j$ lies in $U$, and we may extend a path from $\vec x_0$ to $\vec x$ by this segment. Hence
\[
f(\vec x+h\vec e_j)-f(\vec x)=\int_0^hF_j(\vec x+s\vec e_j)\,\dd s,\qquad\text{so}\qquad\frac{f(\vec x+h\vec e_j)-f(\vec x)}h\to F_j(\vec x)
\]
as $h\to0$, by continuity of $F_j$ (mean value theorem for integrals). Thus $\partial_jf=F_j$ exists and is continuous. By Theorem~\ref{thm:c1diff}, $f$ is $C^1$ with $\nabla f=\vec F$.
\end{proof}

\begin{intuition}
Conservative means that the work depends only on where you start and where you end, as with the gravitational field: climbing a hill costs the same energy by any route. The potential $f$ is like the height function, and the line integral is the change in height. For a non-conservative field (wind along a circular storm), going around a loop can give work for free.
\end{intuition}

\begin{theorem}[Necessary condition]\label{thm:necessary}
If $\vec F$ is $C^1$ and conservative on $U$, then $\curl\vec F=\vec 0$ on $U$ (in $\R^2$: $\partial_xQ=\partial_yP$).
\end{theorem}

\begin{proof}
$\vec F=\nabla f$ with $f$ of class $C^2$ (as $\vec F$ is $C^1$), so $\curl\vec F=\curl\nabla f=\vec 0$ by Theorem~\ref{thm:identities}(1).
\end{proof}

The converse is true on star-shaped regions (Theorem~\ref{thm:poincare}), and more generally on simply connected regions. But it is false in general, as the next section shows.

\begin{example}[Finding a potential and computing a line integral]
The field $\vec F=(2xy+z^2,\ x^2,\ 2xz)$ has $\curl\vec F=\vec 0$ on $\R^3$ and potential $f=x^2y+xz^2$ (Chapter~\ref{ch:fields}). The work along any curve from $(0,0,0)$ to $(1,2,3)$ is $f(1,2,3)-f(0,0,0)=1\cdot2+1\cdot9=11$.
\end{example}

%====================================================================
\section{The vortex field and the winding number}
%====================================================================

Consider on $U=\R^2\setminus\{(0,0)\}$ the field
\[
\vec G(x,y)=\Bigl(\frac{-y}{x^2+y^2},\ \frac{x}{x^2+y^2}\Bigr)=\frac1r\,\ehat_\theta .
\]
We computed in Chapter~\ref{ch:fields} that $\curl\vec G=\vec 0$. Nevertheless:

\begin{theorem}[The vortex is not conservative]\label{thm:vortex}
For the unit circle $C$, traversed counterclockwise, $\oint_C\vec G\cdot\dd\vec r=2\pi\neq0$. Hence $\vec G$ is not conservative on $\R^2\setminus\{\vec 0\}$, although $\curl\vec G=\vec 0$ there.
\end{theorem}

\begin{proof}
With $\vec r(t)=(\cos t,\sin t)$ we have $\vec G(\vec r(t))=(-\sin t,\cos t)=\vec r\,'(t)$ (since $x^2+y^2=1$), so $\vec G\cdot\vec r\,'=1$ and $\oint_C\vec G\cdot\dd\vec r=\int_0^{2\pi}\dd t=2\pi$. By Theorem~\ref{thm:pathind}(c), $\vec G$ is not conservative.
\end{proof}

Why does this not contradict the Poincar\'e lemma? Because $\R^2\setminus\{\vec 0\}$ has a \emph{hole} and is not star-shaped. The field $\vec G$ \emph{is} locally a gradient: it is $\nabla\theta$, where $\theta=\arg(x,y)$ is the polar angle. But the angle cannot be defined continuously on the whole punctured plane: going once around the origin, $\theta$ increases by $2\pi$ and does not return to its starting value.

\begin{theorem}[Winding number]\label{thm:winding}
Let $C$ be a closed piecewise $C^1$ curve in $\R^2\setminus\{\vec 0\}$. Then
\[
\oint_C\vec G\cdot\dd\vec r=2\pi\,w(C),
\]
where the integer $w(C)$ is the number of times $C$ winds counterclockwise around the origin (the \emph{winding number}).
\end{theorem}

\begin{proof}
Write the curve in polar form $\vec r(t)=\rho(t)\bigl(\cos\theta(t),\sin\theta(t)\bigr)$, where $\rho=\norm{\vec r}>0$ and $\theta(t)$ is a continuous choice of the polar angle along the curve (it exists because the angle can be chosen continuously along any path, by following it step by step; this is the standard ``lifting'' property of the angle). Then $\vec r\,'=\rho'\ehat_r+\rho\theta'\ehat_\theta$ (Theorem~\ref{thm:polaracc}) and $\vec G=\rho^{-1}\ehat_\theta$, so $\vec G\cdot\vec r\,'=\theta'(t)$. Therefore
\[
\oint_C\vec G\cdot\dd\vec r=\int_a^b\theta'(t)\,\dd t=\theta(b)-\theta(a).
\]
Since $\vec r(b)=\vec r(a)$, the angles $\theta(b)$ and $\theta(a)$ differ by a multiple of $2\pi$; the multiple is by definition the winding number.
\end{proof}

\begin{figure}[ht]
\centering
\begin{tikzpicture}[scale=1.15]
  \foreach \xs/\what in {0/{$w=1$},4.2/{$w=0$},8.4/{$w=2$}}{
    \begin{scope}[xshift=\xs cm]
      \foreach \x in {-1.2,-0.6,0,0.6,1.2}{ \foreach \y in {-1.2,-0.6,0,0.6,1.2}{
        \pgfmathparse{(\x==0 && \y==0) ? 0 : 1}
        \ifnum\pgfmathresult>0
          \pgfmathsetmacro{\nn}{sqrt(\x*\x+\y*\y)}
          \draw[-{Stealth[length=1.3mm]},mid!60] (\x,\y)--++({0.3*(-\y)/\nn},{0.3*\x/\nn});
        \fi
      }}
      \node[pt] at (0,0) {};
      \node[font=\small] at (0,-1.75) {\what};
    \end{scope}
  }
  % w=1: circle around origin
  \draw[curvestyle,very thick,postaction={decorate,decoration={markings,mark=at position 0.15 with {\arrow{Stealth}}}}] (0,0) circle (0.95);
  % w=0: circle not around origin
  \draw[accent2,very thick,postaction={decorate,decoration={markings,mark=at position 0.15 with {\arrow{Stealth}}}}] (4.2+0.65,0.55) circle (0.45);
  % w=2: double loop spiral
  \draw[accent3,very thick] plot[domain=0:720,samples=160,smooth]({8.4+(0.45+0.0007*\x)*cos(\x)},{(0.45+0.0007*\x)*sin(\x)});
\end{tikzpicture}
\caption{The vortex field $\vec G=\frac1r\ehat_\theta$ (grey arrows; the origin is removed) and three closed curves. The circulation $\oint\vec G\cdot\dd\vec r$ equals $2\pi w$, where $w$ is the number of counterclockwise turns around the origin: $2\pi$, $0$ and $4\pi$.}
\label{fig:winding}
\end{figure}

\begin{pitfall}
``$\curl\vec F=\vec0$ implies conservative'' is true only on regions without holes (star-shaped, or simply connected). The domain is part of the hypothesis. This is a very common interview question, with the vortex $\vec G$ as the standard answer.
\end{pitfall}

%====================================================================
\section{Parametrized surfaces and surface area}
%====================================================================

A \emph{parametrized surface} is a map $\vec r:D\to\R^3$, $\vec r(u,v)=(x(u,v),y(u,v),z(u,v))$, from a region $D\subseteq\R^2$, of class $C^1$. The vectors
\[
\vec r_u=\pd{\vec r}u,\qquad\vec r_v=\pd{\vec r}v
\]
are tangent to the coordinate curves on the surface. The surface is \emph{regular} if $\vec r_u\times\vec r_v\ne\vec 0$ everywhere, so that $\vec r_u,\vec r_v$ span the tangent plane. The vector $\vec N=\vec r_u\times\vec r_v$ is normal to the surface, and $\hat{\vec n}=\vec N/\norm{\vec N}$ is a unit normal.

\begin{figure}[ht]
\centering
\begin{tikzpicture}[scale=1.25]
  % left: domain
  \draw[axisline,mid] (-0.3,0)--(2.8,0) node[right]{$u$};
  \draw[axisline,mid] (0,-0.3)--(0,2.6) node[above]{$v$};
  \filldraw[fill=accent!12,draw=accent,thick] (0.3,0.3) rectangle (2.4,2.1);
  \filldraw[fill=accent2!35,draw=accent2,thick] (1.2,0.9) rectangle (1.7,1.4);
  \node[font=\scriptsize,accent2] at (1.45,0.62) {$\dd u\,\dd v$};
  \draw[-{Stealth[length=3mm]},very thick,mid] (3.1,1.3)--(4.2,1.3) node[midway,above]{$\vec r$};
  % right: surface patch
  \begin{scope}[xshift=4.6cm]
    \path[fill=accent!12,draw=accent,thick] (0,0.1) to[out=-10,in=190] (3.3,0.4) to[out=80,in=-95] (3.7,2.5) to[out=170,in=5] (0.5,2.2) to[out=-100,in=95] cycle;
    \path[fill=accent!12] (0,0.1) to[out=-10,in=190] (3.3,0.4) to[out=80,in=-95] (3.7,2.5) to[out=170,in=5] (0.5,2.2) to[out=-100,in=95] cycle;
    \draw[accent,thick] (0,0.1) to[out=-10,in=190] (3.3,0.4) to[out=80,in=-95] (3.7,2.5) to[out=170,in=5] (0.5,2.2) to[out=-100,in=95] cycle;
    \coordinate (P) at (1.6,1.1);
    \node[pt] at (P) {};
    \draw[vecstyle,accent2] (P)--++(1.0,0.2) node[right,font=\small]{$\vec r_u$};
    \draw[vecstyle,accent3] (P)--++(0.15,0.95) node[above,font=\small]{$\vec r_v$};
    \filldraw[fill=accent2!35,draw=accent2,thick,opacity=0.8] (P)--++(0.6,0.12)--++(0.09,0.57)--++(-0.6,-0.12)--cycle;
    \draw[vecstyle,accent4] (P)--++(-0.5,0.85) node[left,font=\small]{$\vec N$};
    \node[accent,font=\small] at (3.0,1.9) {$S$};
  \end{scope}
\end{tikzpicture}
\caption{A parametrization maps a small rectangle $\dd u\,\dd v$ to a nearly flat patch on the surface. This patch is almost the parallelogram spanned by $\vec r_u\dd u$ and $\vec r_v\dd v$, with area $\norm{\vec r_u\times\vec r_v}\dd u\,\dd v$.}
\label{fig:surfacepatch}
\end{figure}

By Theorem~\ref{thm:crossprops}(c), the parallelogram spanned by $\vec r_u\,\dd u$ and $\vec r_v\,\dd v$ has area $\norm{\vec r_u\times\vec r_v}\,\dd u\,\dd v$. This leads to the definition
\[
\Area(S)=\iint_D\norm{\vec r_u\times\vec r_v}\,\dd u\,\dd v,\qquad\dd S=\norm{\vec r_u\times\vec r_v}\,\dd u\,\dd v .
\]

\begin{proposition}[Independence of parametrization]\label{prop:surfparam}
Let $\vec\varphi:D'\to D$, $(s,t)\mapsto(u,v)$ be a $C^1$ bijection with $\det\Jac\vec\varphi\neq0$ and let $\tilde{\vec r}=\vec r\circ\vec\varphi$. Then
\[
\tilde{\vec r}_s\times\tilde{\vec r}_t=\det\Jac\vec\varphi\cdot(\vec r_u\times\vec r_v)\circ\vec\varphi,
\]
and the surface area computed with $\tilde{\vec r}$ equals the area computed with $\vec r$.
\end{proposition}

\begin{proof}
By the chain rule, $\tilde{\vec r}_s=\vec r_uu_s+\vec r_vv_s$ and $\tilde{\vec r}_t=\vec r_uu_t+\vec r_vv_t$. Using bilinearity of the cross product and $\vec r_u\times\vec r_u=\vec r_v\times\vec r_v=\vec 0$, $\vec r_v\times\vec r_u=-\vec r_u\times\vec r_v$:
\[
\tilde{\vec r}_s\times\tilde{\vec r}_t=(u_sv_t-v_su_t)\,\vec r_u\times\vec r_v=\det\begin{pmatrix}u_s&u_t\\v_s&v_t\end{pmatrix}\,\vec r_u\times\vec r_v .
\]
Taking norms and applying the change-of-variables formula (Theorem~\ref{thm:cov}),
\[
\iint_{D'}\norm{\tilde{\vec r}_s\times\tilde{\vec r}_t}\dd s\,\dd t=\iint_{D'}\norm{\vec r_u\times\vec r_v}\,\abs{\det\Jac\vec\varphi}\,\dd s\,\dd t=\iint_D\norm{\vec r_u\times\vec r_v}\dd u\,\dd v .\qedhere
\]
\end{proof}

Note also that the unit normal direction $\vec r_u\times\vec r_v$ is preserved when $\det\Jac\vec\varphi>0$ and reversed when $\det\Jac\vec\varphi<0$.

\begin{example}[Standard surface elements]
\begin{enumerate}[label=(\alph*)]
  \item \emph{Graph} $z=f(x,y)$: $\vec r(x,y)=(x,y,f)$, $\vec r_x=(1,0,f_x)$, $\vec r_y=(0,1,f_y)$, so $\vec N=\vec r_x\times\vec r_y=(-f_x,-f_y,1)$ (pointing upward) and
  \[
  \dd S=\sqrt{1+f_x^2+f_y^2}\,\dd x\,\dd y .
  \]
  \item \emph{Sphere} of radius $R$: $\vec r(\varphi,\theta)=R(\sin\varphi\cos\theta,\sin\varphi\sin\theta,\cos\varphi)$. Then $\vec r_\varphi\times\vec r_\theta=R^2\sin\varphi\,\ehat_\rho$ (outward), $\dd S=R^2\sin\varphi\,\dd\varphi\,\dd\theta$, and the area is $4\pi R^2$.
  \item \emph{Cylinder} $x^2+y^2=R^2$, $0\le z\le h$: $\vec r(\theta,z)=(R\cos\theta,R\sin\theta,z)$, $\vec r_\theta\times\vec r_z=(R\cos\theta,R\sin\theta,0)$ (outward), $\dd S=R\,\dd\theta\,\dd z$, area $2\pi Rh$.
  \item \emph{Paraboloid} $z=x^2+y^2$, $x^2+y^2\le1$: $\dd S=\sqrt{1+4x^2+4y^2}\,\dd A$. In polar coordinates the area is $\int_0^{2\pi}\!\int_0^1\sqrt{1+4r^2}\,r\,\dd r\,\dd\theta=2\pi\cdot\frac1{12}\bigl[(1+4r^2)^{3/2}\bigr]_0^1=\frac\pi6\bigl(5\sqrt5-1\bigr)$.
\end{enumerate}
\end{example}

%====================================================================
\section{Surface integrals and flux}
%====================================================================

\begin{definition}[Surface integrals]
For a scalar function $f$ and a parametrized surface $S$,
\[
\iint_Sf\,\dd S=\iint_Df\bigl(\vec r(u,v)\bigr)\,\norm{\vec r_u\times\vec r_v}\,\dd u\,\dd v .
\]
For a vector field $\vec F$ and an \emph{oriented} surface (a surface with a chosen unit normal $\hat{\vec n}$, here $\hat{\vec n}=\vec r_u\times\vec r_v/\norm{\vec r_u\times\vec r_v}$), the \emph{flux} of $\vec F$ through $S$ is
\[
\iint_S\vec F\cdot\dd\vec S=\iint_S\vec F\cdot\hat{\vec n}\,\dd S=\iint_D\vec F\bigl(\vec r(u,v)\bigr)\cdot\bigl(\vec r_u\times\vec r_v\bigr)\,\dd u\,\dd v .
\]
\end{definition}

The flux is the volume of fluid that crosses $S$ per unit time (in the direction of $\hat{\vec n}$) if $\vec F$ is the velocity of the fluid. By Proposition~\ref{prop:surfparam}, the scalar surface integral does not depend on the parametrization, and the flux depends only on the surface and its orientation (it changes sign if the orientation is reversed).

\begin{theorem}[Flux through a graph]\label{thm:fluxgraph}
For the graph $z=f(x,y)$ over a region $D$, oriented upward, and $\vec F=(P,Q,R)$,
\[
\iint_S\vec F\cdot\dd\vec S=\iint_D\bigl(-P\,f_x-Q\,f_y+R\bigr)\,\dd x\,\dd y,
\]
with $P,Q,R$ evaluated at $(x,y,f(x,y))$.
\end{theorem}

\begin{proof}
Here $\vec r_x\times\vec r_y=(-f_x,-f_y,1)$ (see the example above); take the dot product with $\vec F$.
\end{proof}

\begin{example}[Flux through a sphere]
Let $\vec F=\vec r/\norm{\vec r}^3$ and let $S$ be the sphere of radius $R$ with outward normal $\ehat_\rho$. On $S$, $\vec F\cdot\hat{\vec n}=R^{-2}$, which is constant, so
\[
\iint_S\vec F\cdot\dd\vec S=\frac1{R^2}\cdot\Area(S)=\frac{4\pi R^2}{R^2}=4\pi .
\]
The flux does not depend on $R$! This is the source of Gauss's law (Chapter~\ref{ch:physics}). For the position field $\vec F=\vec r$ we get $R\cdot4\pi R^2=4\pi R^3$.
\end{example}

\begin{example}[Flux through a cylinder wall]
For $\vec F=(x,y,0)$ and the cylinder wall of the example above (radius $R$, height $h$, outward normal): $\vec F\cdot(\vec r_\theta\times\vec r_z)=R^2\cos^2\theta+R^2\sin^2\theta=R^2$, so the flux is $\int_0^h\!\int_0^{2\pi}R^2\,\dd\theta\,\dd z=2\pi R^2h$. (Compare with $\divg\vec F=2$ times the volume $\pi R^2h$, which we explain in Chapter~\ref{ch:theorems}.)
\end{example}

\subsection{Orientability}

A smooth surface is \emph{orientable} if one can choose a unit normal $\hat{\vec n}$ that varies continuously over the whole surface. Spheres, discs, graphs and tori are orientable. A flux integral requires an orientation.

\begin{figure}[ht]
\centering
\begin{tikzpicture}
\begin{axis}[width=9cm,view={60}{35},hide axis,colormap={b}{color=(softblue) color=(accent)},z buffer=sort,
  domain=0:360,y domain=-0.45:0.45,samples=60,samples y=6]
\addplot3[surf,shader=faceted interp,faceted color=ink!25]
  ({(1+y*cos(x/2))*cos(x)},{(1+y*cos(x/2))*sin(x)},{y*sin(x/2)});
\end{axis}
\end{tikzpicture}
\caption{The M\"obius strip. A unit normal moved once around the centre line returns pointing in the opposite direction, so the strip has only one side and is \emph{not orientable}. Flux through it is not defined.}
\label{fig:mobius}
\end{figure}

%====================================================================
\begin{summarybox}
\begin{itemize}
  \item $\int_Cf\,\dd s=\int f(\vec r)\norm{\vec r\,'}\dd t$ (no orientation needed). $\int_C\vec F\cdot\dd\vec r=\int\vec F(\vec r)\cdot\vec r\,'\dd t$ (orientation changes the sign).
  \item $\int_C\nabla f\cdot\dd\vec r=f(\text{end})-f(\text{start})$. For continuous $\vec F$ on a connected open set: conservative $\Leftrightarrow$ path independent $\Leftrightarrow$ all closed integrals vanish.
  \item $\curl\vec F=\vec 0$ is necessary for being conservative, and sufficient on star-shaped or simply connected regions. The vortex $\frac1r\ehat_\theta$ has $\curl=\vec 0$ but circulation $2\pi w$.
  \item Surface: $\dd S=\norm{\vec r_u\times\vec r_v}\dd u\,\dd v$; flux $=\iint\vec F\cdot(\vec r_u\times\vec r_v)\dd u\,\dd v$; graph: $\dd S=\sqrt{1+f_x^2+f_y^2}\dd A$.
  \item Reparametrization multiplies $\vec r_u\times\vec r_v$ by $\det\Jac$; area is invariant, flux changes sign if orientation reverses.
  \item The flux of $\vec r/r^3$ through any sphere around the origin is $4\pi$.
\end{itemize}
\end{summarybox}

\begin{interview}
\begin{enumerate}
  \item \emph{What is the difference between $\int f\,\dd s$ and $\int\vec F\cdot\dd\vec r$?} The first is orientation-independent; the second uses the tangential component and changes sign with orientation.
  \item \emph{State and prove the fundamental theorem for line integrals.} $\int_C\nabla f\cdot\dd\vec r=f(B)-f(A)$; the integrand is $\frac{\dd}{\dd t}f(\vec r(t))$.
  \item \emph{When is a field conservative?} Iff closed-loop integrals vanish iff path independent iff a potential exists. $\curl=0$ is enough on simply connected domains.
  \item \emph{Can a field with zero curl be non-conservative?} Yes: $(-y,x)/(x^2+y^2)$ on the punctured plane; circulation $2\pi$.
  \item \emph{Why is $\dd S=\norm{\vec r_u\times\vec r_v}\dd u\dd v$?} Area of the parallelogram spanned by $\vec r_u\dd u$ and $\vec r_v\dd v$.
  \item \emph{Why does flux need an orientation?} Flux counts the direction of crossing; the choice of normal decides the sign.
  \item \emph{What is the flux of $\vec r/r^3$ through a sphere?} $4\pi$, independent of the radius.
\end{enumerate}
\end{interview}

\subsection*{Exercises}
\begin{exercise}
Compute $\int_C(x\,\dd y-y\,\dd x)$ along the ellipse $x=a\cos t$, $y=b\sin t$, $0\le t\le2\pi$. (Answer: $2\pi ab$, twice the area of the ellipse.)
\end{exercise}
\begin{exercise}
Is $\vec F=(e^x\sin y+y,\ e^x\cos y+x)$ conservative on $\R^2$? Find a potential. (Answer: $f=e^x\sin y+xy$.)
\end{exercise}
\begin{exercise}
Compute the area of the part of the sphere $x^2+y^2+z^2=R^2$ between the planes $z=a$ and $z=b$ ($-R\le a<b\le R$). (Answer: $2\pi R(b-a)$; Archimedes' hat-box theorem.)
\end{exercise}
\begin{exercise}
Find the flux of $\vec F=(x,y,z)$ through the paraboloid $z=1-x^2-y^2$, $z\ge0$, oriented upward, using Theorem~\ref{thm:fluxgraph}. (Answer: $\frac{3\pi}2$.)
\end{exercise}
\begin{exercise}
Show that the surface area of a surface of revolution $\vec r(\theta,x)=(x,g(x)\cos\theta,g(x)\sin\theta)$ is $2\pi\int_a^bg\sqrt{1+g'^2}\,\dd x$.
\end{exercise}
\begin{exercise}
Let $\vec F=\vec G+\nabla h$ where $\vec G$ is the vortex and $h$ is a $C^1$ function on the punctured plane. Show $\oint_C\vec F\cdot\dd\vec r=2\pi w(C)$ for every closed curve $C$.
\end{exercise}
